\section{Introducción}

En 1948, John von Neumann, con ayuda de Stanislaw Ulam, propuso un Autómata Celular (CA, \textit{Cellular Automaton}) como un modelo matemático con la capacidad de autoreproducción universal, el cual desarrollo en su libro \textit{Theory of Self-Reproducing Automata} \cite{vonNeumann1948}. Hoy en día representa un modelo fundamental para el estudio de sistemas complejos, computación distribuida y procesos autoorganizados. 

Diversos autores, como Stephen Wolfram \cite{wolfram1982} y Erick W. Weisstein \cite{weissteinCA} coinciden en caracterizar a los CA como sistemas discretos compuestos por células dispuestos sobre una retícula\footnote{Conjunto de líneas o elementos de estructura lineal que se cruzan o entrelazan estrechamente \cite{fundacionBBVALattice}} (\textit{lattice}). Cada celda puede adoptar un valor de un conjunto finito de estados, y su actualización ocurre en pasos discretos de tiempo, determinada por una función de transición que depende únicamente de su vecindad local.

Aunque existen variaciones en la terminología (celda, célula, sitio, color), la estructura formal es ampliamente aceptada y utilizada en campos como la física, biología computacional, procesamiento de imágenes, teoría de la computación, entre otros. Más recientemente, se han integrado técnicas de aprendizaje supervisado, en particular Redes Neuronales Artificiales (ANN), dando lugar a los Autómatas Celulares Neuronales (NCA, \textit{Neural Cellular Automata}).

El presente estado del arte revisa estas propuestas híbridas, enfatizando sus aplicaciones en morfogénesis artificial, clasificación de patrones, generación de texturas dinámicas y modelado en mallas tridimensionales. Además, se identifican retos en términos de escalabilidad, generalización y control de dinámicas emergentes, aspectos que motivan el desarrollo de nuevas variantes de NCA.
